%% kapa: exam=1 round=25 year=2014 period=1 form=A subject=civil_law subject_ko=민법 seq=24 number=24
%% source: Q-Net 감정평가사 기출문제 (cst00309) artlSeq=5101207 file=extracted/1008/5101207/2014년 제25회 감정평가사 1차시험 기출문제/1교시(A형).hwp
%% conversion: kordoc 4.24.51 -> tools/gen1.py (auto); review=proofread-sonnet; flags=none
\begin{question}{24}
법률행위의 무효와 취소에 관한 설명으로 옳지 않은 것은? (다툼이 있으면 판례에 의함)
\begin{choices}
  \item 당사자가 무효임을 알고 추인한 때에는 새로운 법률행위를 한 것으로 본다.
  \item 취소할 수 있는 법률행위의 취소원인이 소멸한 후에 취소권자의 상대방이 이행을 청구한 경우에는 다른 사정이 없는 한, 법률상 당연히 추인이 있었던 것으로 본다.
  \item 취소할 수 있는 법률행위의 상대방이 확정된 경우에는, 그 취소는 그 상대방에 대한 의사표시로 하여야 한다.
  \item 제한능력자가 제한능력을 이유로 법률행위를 취소한 경우에는, 그 행위로 인하여 받은 이익이 현존하는 한도에서 상환책임이 있다.
  \item 취소할 수 있는 법률행위가 적법하게 취소되어 무효가 되었더라도 그 무효행위가 추인의 요건을 갖춘 경우에는 이를 다시 추인할 수 있다.
\end{choices}
\end{question}
\answer{2}
